% =====================================================================================
% REVISIÓN 3  -  PROYECTOS DE SISTEMA INFORMÁTICO (página web, sistema web, API, escritorio, app móvil)
% Para usarla pon \TipoProyecto{sistema} en configuracion/datos.tex y deja activa esta línea en main.tex.
% Escribe tu texto FUERA de los recuadros \begin{guia} ... \end{guia}.
% =====================================================================================

\section{Desarrollo del sistema} \label{sec:desarrollo_sistema} % audit:req=desarrollo_sistema

\begin{guia}[Guía: Desarrollo del sistema (Revisión 3, proyectos de sistema)]
\textbf{¿Para qué sirve?} Documenta cómo se construyó la solución: arquitectura, diseño y resultado final. Debe permitir
que otra persona entienda el sistema y pueda darle mantenimiento.

\textbf{En el texto propio de este capítulo (antes de las subsecciones) responde:}
\begin{itemize}[nosep,leftmargin=*]
  \item ¿Cuál es la arquitectura del sistema (capas, cliente-servidor, microservicios, nube)? Incluye un diagrama.
  \item ¿Qué tecnologías, versiones y librerías se usaron y para qué?
  \item ¿Cómo está organizado el código (módulos, carpetas, patrones)? Muestra solo fragmentos relevantes y explícalos.
  \item ¿Qué pruebas se hicieron (funcionales, de integración, de usabilidad) y con qué resultado?
  \item ¿Cómo se despliega o instala? ¿dónde está el repositorio?
  \item ¿Qué problemas técnicos surgieron y cómo se resolvieron?
\end{itemize}

\textbf{Extensión mínima:} 300 palabras de texto propio de este capítulo, además de las subsecciones que siguen.

\textbf{Reglas para todas las figuras:} cada diagrama o captura lleva \verb|\caption|, \verb|\label|, se cita en el texto con
\verb|Figura~\ref{fig:...}| (nunca escribas ``Figura 3'' a mano) y debe ser legible al imprimirse.

\textbf{Cómo se revisa:} coherencia con la propuesta y los requisitos, suficiencia técnica y que las subsecciones
obligatorias estén completas.
\end{guia}

% >>> ESCRIBE AQUÍ la arquitectura, tecnologías, implementación y pruebas <<<


\subsection{Diagramas de casos de uso} \label{sec:casos_uso} % audit:req=casos_uso

\begin{guia}[Guía: Diagramas de casos de uso]
\textbf{Debe incluir:} un diagrama general con los actores y casos de uso (relaciones \emph{include} y \emph{extend} cuando
apliquen) y la \textbf{especificación} de los casos de uso principales en una tabla (nombre, actor, precondición, flujo
normal, flujos alternos, postcondición).

\textbf{Debe responder:} ¿quiénes usan el sistema? ¿qué puede hacer cada actor? ¿cada caso de uso corresponde a un
requisito y a un objetivo específico?

\textbf{Mínimo:} 1 diagrama (figura) con explicación. Herramientas: draw.io, StarUML, PlantUML, Visual Paradigm.
\end{guia}

% >>> ESCRIBE AQUÍ los diagramas de casos de uso y su explicación <<<


\subsection{Diagramas de base de datos} \label{sec:diagramas_bd} % audit:req=diagrama_bd

\begin{guia}[Guía: Diagramas de base de datos]
\textbf{Debe incluir:} el diagrama entidad-relación (o el modelo relacional / de documentos) y el
\textbf{diccionario de datos} (tabla por entidad: campo, tipo, restricciones, descripción).

\textbf{Debe responder:} ¿qué entidades y relaciones hay y por qué? ¿qué gestor de base de datos se usó y por qué? ¿está
normalizado? ¿cómo se protege la información (respaldos, permisos, datos personales)?

\textbf{Mínimo:} 1 diagrama (figura) con explicación.
\end{guia}

% >>> ESCRIBE AQUÍ los diagramas de base de datos y el diccionario de datos <<<


\subsection{Diagramas de actividades} \label{sec:diagramas_actividades} % audit:req=diagrama_actividades

\begin{guia}[Guía: Diagramas de actividades]
\textbf{Debe incluir:} el flujo de los procesos clave del sistema (con carriles o \emph{swimlanes} por actor cuando hay
varios participantes), con inicio, decisiones, acciones y fin.

\textbf{Debe responder:} ¿cuál es el flujo de cada proceso importante? ¿qué pasa en los casos de error? ¿corresponde a
algún caso de uso?

\textbf{Mínimo:} 1 diagrama (figura) con explicación.
\end{guia}

% >>> ESCRIBE AQUÍ los diagramas de actividades <<<


\subsection{Diagramas de secuencia} \label{sec:diagramas_secuencia} % audit:req=diagrama_secuencia

\begin{guia}[Guía: Diagramas de secuencia]
\textbf{Debe incluir:} la interacción entre los componentes (usuario, interfaz, servidor, base de datos, servicios
externos) para los casos de uso críticos, con los mensajes en orden temporal.

\textbf{Debe responder:} ¿qué componentes intervienen y en qué orden se comunican? ¿qué mensajes se envían y qué se
responde? ¿dónde se validan los datos?

\textbf{Mínimo:} 1 diagrama (figura) con explicación.
\end{guia}

% >>> ESCRIBE AQUÍ los diagramas de secuencia <<<


\subsection{Resultado final: pantallas del sistema} \label{sec:pantallas} % audit:req=pantallas

\begin{guia}[Guía: Resultado final -- pantallas del sistema con su explicación]
\textbf{Debe incluir:} capturas de las pantallas principales del sistema funcionando, cada una con su explicación.

\textbf{Para cada pantalla responde:} ¿para qué sirve? ¿quién la usa? ¿qué casos de uso y objetivos específicos
atiende? ¿qué datos muestra o captura y cómo se valida?

\textbf{Reglas:} figuras numeradas con \verb|\caption| y \verb|\label|; sin datos personales ni confidenciales reales
(usa datos de prueba); capturas legibles. Si hay versión móvil o de escritorio, inclúyela.

\textbf{Mínimo:} 150 palabras de explicación y al menos una figura.
\end{guia}

% >>> ESCRIBE AQUÍ las pantallas (figuras) y su explicación <<<
